% Fragmento de entrega (sin preambulo; se compila desde main.tex).
\section*{Problema 1 -- Teoría de equidad e imposibilidad}

En lo que sigue $Y=1$ significa que el cliente incumplió, $Y=0$ que pagó, y
$\hat{Y}=1$ que el banco rechazó el crédito. Para un grupo $G=g$ escribimos
$p_g=P(Y=1\mid G=g)$, $PPV_g=P(Y=1\mid \hat{Y}=1,G=g)$,
$FPR_g=P(\hat{Y}=1\mid Y=0,G=g)$ y $FNR_g=P(\hat{Y}=0\mid Y=1,G=g)$.

\subsection*{(a) La identidad del FPR}

Escribimos $P(\hat{Y}=1,Y=0)$ de dos formas. Condicionando sobre $Y$,
\[
P(\hat{Y}=1,Y=0)=P(\hat{Y}=1\mid Y=0)\,P(Y=0)=FPR\,(1-p).
\]
Condicionando sobre $\hat{Y}$,
\[
P(\hat{Y}=1,Y=0)=P(Y=0\mid \hat{Y}=1)\,P(\hat{Y}=1)=(1-PPV)\,P(\hat{Y}=1).
\]
Análogamente para los incumplidores rechazados,
\begin{align*}
P(\hat{Y}=1,Y=1)&=P(\hat{Y}=1\mid Y=1)\,P(Y=1)=(1-FNR)\,p\\
&=P(Y=1\mid \hat{Y}=1)\,P(\hat{Y}=1)=PPV\,P(\hat{Y}=1).
\end{align*}
Dividiendo la primera expresión entre la segunda se cancela $P(\hat{Y}=1)$
(suponiendo que se rechaza a alguien) y se obtiene
\[
FPR=\frac{(1-PPV)\,P(\hat{Y}=1)}{(1-p)}
    \quad\text{y}\quad
    (1-FNR)\,p = PPV\,P(\hat{Y}=1),
\]
de modo que
\[
\;FPR=\frac{p}{1-p}\cdot\frac{1-PPV}{PPV}\cdot(1-FNR)\;
\]
Forma equivalente por conteos: con $P$ positivos reales y $N$ negativos reales,
\[
FPR=\frac{FP}{N}=\frac{FP}{TP}\cdot\frac{TP}{P}\cdot\frac{P}{N}
   =\frac{1-PPV}{PPV}\,(1-FNR)\,\frac{p}{1-p}.
\]

La identidad es puramente algebraica: no es una coincidencia empírica, sino
una consecuencia de usar la misma matriz de confusión. En la práctica sirve
para verificar que una tabla por grupo es internamente consistente.

\subsection*{(b) Igualar FPR, PPV y FNR es imposible con tasas base distintas}

El enunciado pide comparar dos grupos $A$ y $B$ con $FPR_A=0.30$ y
$FPR_B=0.175$. Si se quisiera igualar además la PPV y el FNR, la identidad del
inciso (a) daría un valor requerido de $PPV_B=7/19\approx 0.368$. El caso
general muestra que eso no puede funcionar cuando las tasas base difieren.

Supongamos que el PPV y el FNR son los mismos en ambos grupos, con
$p_A\neq p_B$. Igualar el FPR exige
\[
\frac{p_A}{1-p_A}\cdot\frac{1-PPV}{PPV}\cdot(1-FNR)
   =\frac{p_B}{1-p_B}\cdot\frac{1-PPV}{PPV}\cdot(1-FNR).
\]
Como $p_A\neq p_B$, los factores $p/(1-p)$ son distintos. Para que la
igualdad pueda cumplirse debe ser
\[
(1-PPV)(1-FNR)=0.
\]
Es decir, hay dos únicas salidas: $PPV=1$, que significa que no hay falsos
positivos (y por tanto $FPR=0$ en ambos grupos), o $FNR=1$, que significa que
nunca se predice positivo, en cuyo caso la PPV queda indefinida porque no hay
rechazos. Por eso la frase del enunciado ``no perfecto'' debe leerse como
``comete falsos positivos'': el teorema no dice que no exista un clasificador
justo en sentido absoluto, sino que un clasificador con algún falso positivo
no puede igualar las tres métricas a la vez con tasas base distintas. Vale
la pena observar por separado el caso $PPV=1$ con $FNR>0$: allí nunca se
rechaza a un cliente que iba a pagar, pero sí se aprueban algunos que
incumplen; es un clasificador selectivo en el rechazo, no perfecto. El
clasificador perfecto es el caso extremo $FPR=FNR=0$ y $PPV=1$, que sí iguala
todas las métricas trivialmente porque no comete ningún error. Fuera de ese
caso límite y de los degenerados, la imposibilidad entra en juego.

\subsection*{(c) Calibración y los errores por grupo}

El enunciado plantea dos supuestos: (P1) la calibración por grupo,
$P(Y=1\mid S=s,G=g)=s$, donde $S$ es el puntaje, y (P2) la ley de
probabilidad total. Entonces
\[
E[S\mid G=g]=E\bigl[E[Y\mid S,G]\mid G=g\bigr]=E[Y\mid G=g]=p_g,
\]
porque $E[Y\mid S,G]=P(Y=1\mid S,G)=S$ bajo (P1). Es decir, un puntaje
calibrado tiene media igual a la tasa base de su grupo: su comprobación
numérica es una prueba de calibración.

El ejemplo del enunciado toma $p_A=11/25=0.44$ y $p_B=13/50=0.26$, y
considera rechazar exactamente a quienes tienen $S=0.8$, cuyo PPV es $0.8$
por calibración. Aplicando la identidad se obtiene
\[
FPR_A=\frac{1}{7}\approx 0.143,\qquad
FNR_A=\frac{3}{11}\approx 0.273,
\]
\[
FPR_B=\frac{1}{37}\approx 0.027,\qquad
FNR_B=\frac{9}{13}\approx 0.692.
\]
Con el criterio FPR sale más perjudicado el grupo $A$, porque se rechaza a
una fracción mayor de quienes no incumplen; con el criterio FNR sale más
perjudicado el grupo $B$, porque se aprueba a una fracción mayor de quienes
sí incumplen. Ninguna de las dos lecturas es la ``correcta'': cada una
prioriza un tipo de error, y con puntajes calibrados y tasas base distintas
no pueden coincidir.

\subsection*{(d) Por qué el umbral óptimo es el mismo para todos}

Dar crédito a un cliente con puntaje calibrado $s$ produce una ganancia
esperada
\[
(1-s)G-sL.
\]
Esta cantidad es positiva si y solo si $(1-s)G>sL$, es decir, si y solo si
\[
s<\frac{G}{G+L}.
\]
Con $L=5G$ el umbral es $\frac{G}{G+5G}=1/6$. El umbral no depende del grupo:
por calibración, $P(Y=1\mid S=s,G=g)=s$ en cada grupo, así que el cálculo
esperado es el mismo para un hombre y para una mujer con el mismo puntaje.
Pero por el inciso (b), aplicar el mismo umbral con tasas base distintas
($p_H\neq p_M$) conduce a FPR y FNR distintas. La ganancia óptima y la
igualdad de errores son objetivos distintos: el primero sugiere un solo
umbral, el segundo exige en general umbrales distintos o intervenir de otra
forma.

\subsection*{(e) Tú decides: definición, umbrales y uso del atributo}

\begin{center}
\fbox{\parbox{0.93\textwidth}{\textbf{Tú decides.}
\textbf{Definición elegida:} igualar el FPR, porque es la diferencia entre
grupos que el bootstrap de 05 distingue del ruido ($0.042289$, IC 95\%
$[0.020183,0.064465]$, sin el cero), mientras que la diferencia de FNR es
compatible con cero ($-0.023922$, IC 95\% $[-0.064162,0.015438]$). En la
cuadrícula $0.05$--$0.60$ con tolerancia $0.01$, el par de máxima ganancia
fuera de muestra es $(t_H,t_M)=(0.19,0.18)$: igual FNR e igual FPR y FNR
eligen el mismo par, así que la elección es robusta a esas variantes.
Paridad demográfica y paridad predictiva se descartan como definición
principal porque dejan diferencias de errores mayores ($1.76$ y $1.26$ puntos
en OOF la primera; $10.96$ y $-7.51$ la segunda), aunque en prueba llegaron a
dar ganancias mayores que la línea base ($+11$ y $+35$ unidades de $G$), ambas
dentro del ruido. El costo de ser justo no se distingue de cero: pérdida de
$7$ unidades de $G$, IC95\% $[-100,116]$. Con $p_H=0.233$ y $p_M=0.213$ la
imposibilidad teórica es leve: con $(0.19,0.18)$ las brechas quedan en FPR
$0.0095$, FNR $0.0078$, PPV $0.016$ y tasa de rechazo $0.014$, así que casi
todas las definiciones se cumplen a la vez. Ojo: los IC95\% de esas brechas
posteriores son anchos (FPR $[-0.012,0.031]$, FNR $[-0.034,0.049]$), así que
su compatibilidad con cero no prueba que las brechas hayan desaparecido, solo
que con esta muestra ya no se distinguen del ruido.

\textbf{Umbrales distintos por grupo:} sí se aceptan como intervención
explícita, con la condición de que se justifiquen y se auditen. El umbral de
hombres es más alto ($0.19$ contra $0.18$) porque su FPR base era mayor: no
es una ventaja arbitraria, es la corrección documentada de la disparidad. No
se usa la prueba para elegirlos: se eligieron con predicciones fuera de
muestra del 70\% y la prueba solo reporta.

\textbf{¿Usar el sexo como entrada del modelo?} No. El contraste de 07
reentrenó el mismo boosting con SEX y obtuvo AUC $0.779162$ (prácticamente el
mismo que $0.779463$ sin SEX), ganancia $2019$ ($+35$, $+1.76\%$), pero
brecha de FPR $0.082428$ e IC sin cero, y brecha de FNR significativa
$-0.053067$. Meter el atributo protegido no elimina las diferencias, las
agrava, y contradice la decisión ya registrada de dejar SEX fuera y usarla
solo para auditoría. Tampoco basta con quitarlo: otras variables predicen
SEX con AUC $0.650929$ en CV y $0.649996$ en prueba (AGE y varias PAY\_*
encabezan la asociación), así que SEX sigue presente en el análisis como
variable de auditoría y, si se interviene, para aplicar el umbral correcto.
El modelo principal se mantiene sin SEX.}}
\end{center}

